\slajdid{06-s022}
\begin{frame}[label=06-s022]{Zaštitni kodovi}
\gusttekst\centering
Na žalost, zbog fizičkih osobina signala, fizike procesa, prilikom obrade, smeštanja i prenosa informacija može doći do grešaka.\\
Bit(i) promeni(e) vrednost sa 0 na 1, ili sa 1 na 0.\\
Da bi se ovo detektovalo i eventualno ispravilo koriste se „zaštitni kodovi“.\par\vspace{2mm}
Bit parnosti\par\vspace{2mm}
\begin{columns}[c,onlytextwidth]
\column{.57\textwidth}\centering\renewcommand{\arraystretch}{1.10}\begin{tabular}{cccc}\toprule\shortstack{Decimalni\\kod} & \shortstack{Binarni\\kod} & Bit parnost & Kompletan kod\\\midrule
0 & 0 0 0 & 0 & 0 0 0 0\\
1 & 0 0 1 & 1 & 0 0 1 1\\
2 & 0 1 0 & 1 & 0 1 0 1\\
3 & 0 1 1 & 0 & 0 1 1 0\\
4 & 1 0 0 & 1 & 1 0 0 1\\
5 & 1 0 1 & 0 & 1 0 1 0\\
6 & 1 1 0 & 0 & 1 1 0 0\\
7 & 1 1 1 & 1 & 1 1 1 1\\
\bottomrule &&&8 od 16\end{tabular}

\column{.40\textwidth}Dodaje se zaštitni bit.\par\vspace{2mm}
Vrednost zaštitnog bita je takva da ukupan broj jedinica u kodu bude paran.\\\centerline{EKSILI}\par\vspace{2mm}
Na prijemu se proverava kod.\par\vspace{2mm}
Može da se detektuje greška nastala na 1, 3, … bita. Neparan broj grešaka.\par\vspace{2mm}
Parna parnost\\Neparna parnost – neparan broj jedinica
\end{columns}
\end{frame}
